%------------------------------------------------------------------------------%
\section{Revisão de Alguns Conceitos de Álgebra}
\label{sec:Revisao_Algebra}

Incluímos aqui uma pequena revisão de alguns conceitos de álgebra que
são necessários para a compreensão do texto.

%------------------------------------------------------------------------------%
\subsection*{Operações Aritméticas}

Para números $a$, $b$, $c$ e $d$, reais ou complexos,
valem as propriedades das operações aritméticas, desde as divisões sejam possíveis.
\begin{align*}
	a(b+c)                                 & = ab + ac                                        &
	\frac{a}{b} + \frac{c}{d}              & = \frac{ad + bc}{bd}                             \\[4mm]
	\frac{a+c}{b}                          & = \frac{a}{b} + \frac{c}{b}                      &
	\dfrac{\;\sfrac{a}{b}\;}{\sfrac{c}{d}} & = \frac{a}{b} \cdot \frac{d}{c} = \frac{ad}{bc}
\end{align*}

%------------------------------------------------------------------------------%
\subsection*{Expoentes e Radicais}

Para números $x$, $y$, $m$ e $n$, reais,
valem as propriedades de expoentes e raízes,
desde as divisões e raízes sejam possíveis.
\begin{align*}
	x^m x^n                     & = x^{m+n}                                 &
	\frac{x^m}{x^n}             & = x^{m-n}                                 \\[5mm]
	(x^m)^n                     & = x^{mn}                                  &
	x^{-n}                      & = \frac{1}{x^n}                           \\[5mm]
	(xy)^n                      & = x^n y^n                                 &
	\left(\frac{x}{y}\right)^n  & = \frac{x^n}{y^n}                         \\[5mm]
	\sqrt[n]{xy\,}              & = \sqrt[n]{x\,}\,\sqrt[n]{y\,}            &
	\sqrt[n]{\frac{\,x\,}{y}\;} & = \frac{\;\sqrt[n]{x\,}\;}{\sqrt[n]{y\,}}
\end{align*}

%------------------------------------------------------------------------------%
\subsection*{Desigualdades e Módulo}

Ao manipular desigualdades vales as seguintes regras
\begin{enumerate}
  \item Se $a<b$ e $b<c$, então $a<c$
  \item Se $a<b$, então $a+c < b+c$
  \item Se $a<b$ e $c>0$, então $ca<cb$
  \item Se $a<b$ e $c<0$, então $ca<cb$
\end{enumerate}
Para $a>0$
\begin{gather*}
	\abs{x} = a \quad\text{significa que}\quad x=a \text{ ou } x=-a     \\
	\abs{x} < a \quad\text{significa que}\quad -a < x < a               \\
	\abs{x} > a \quad\text{significa que}\quad x < -a \text{ ou } a < a
\end{gather*}
Note que os complexos não são ordenados, isso é,
não existem as relações $>$ ou $<$ entre números complexos.

%------------------------------------------------------------------------------%
